\documentclass[11pt]{article}
\usepackage[margin=1in]{geometry}
\usepackage{amsmath,amssymb,amsthm,mathtools}
\usepackage{graphicx}
\usepackage{microtype}
\usepackage[hidelinks]{hyperref}

\newcommand{\C}{\mathbb C}
\newcommand{\Pcone}{\mathbb P}
\newcommand{\Gset}{\Gamma}
\newcommand{\tr}{\operatorname{tr}}
\newcommand{\artanh}{\operatorname{artanh}}
\newtheorem{theorem}{Theorem}
\newtheorem{corollary}{Corollary}
\newtheorem{proposition}{Proposition}

\title{Duality and non-logarithmic nonuniqueness\\
for a multivariable alternative-mean equation}
\author{Candidate partial result for arXiv:2605.16876}
\date{11 August 2026}

\begin{document}
\maketitle

\begin{center}
\fbox{\parbox{0.91\textwidth}{
\textbf{Status: candidate substantial partial result, likely valid.}
Global uniqueness is invariant under \(g(x)\mapsto-g(x^{-1})\), giving the
new unique endpoint \(g(x)=1-x^{-1}\).  Conversely, an explicit
non-logarithmic operator-monotone \(g\) has a certified \(2\times2\)
three-input example with two solutions.  The full classification remains
open.}}
\end{center}

\section{Source and open problem}

The source is Sejong Kim and Vatsalkumar N.~Mer,
\emph{A multivariable mean equation arising from the spectral geometric
mean}, arXiv:2605.16876 (2026) \cite{KM}.  Let \(\Pcone_n\) be the positive
definite \(n\times n\) matrices, let \(\#\) denote the metric geometric mean,
and let
\[
 \mathcal L=\{g:(0,\infty)\to\mathbb R:g\text{ is operator monotone},
 \ g(1)=0,\ g'(1)=1\}.
\]
For a positive probability vector \(\omega=(w_1,\ldots,w_m)\) and
\(\mathbb A=(A_1,\ldots,A_m)\in\Pcone_n^m\), define
\begin{equation}\label{eq:gamma}
 \Gset_g(\omega;\mathbb A)=
 \left\{X\in\Pcone_n:\sum_{i=1}^m w_i g(A_i\#X^{-1})=0\right\}.
\end{equation}
The open problem asks for necessary and sufficient conditions on \(g\in
\mathcal L\) under which this solution is unique.  The source records
uniqueness for \(g(x)=x-1\), where the solution is the Wasserstein mean, and
nonuniqueness for \(g(x)=\log x\).

\begin{figure}[ht]
\centering
\includegraphics[width=\textwidth]{figures/open_problem_crop.png}
\caption{Equation (4.21) and the open classification problem on page 21 of
the source PDF.}
\label{fig:problem}
\end{figure}

\section{Proof intuition}

Inversion supplies an exact symmetry of the problem: replacing \(g\) by
\(-g(1/x)\), inverting every datum, and inverting the unknown leaves the
equation unchanged.  This immediately transfers the Wasserstein uniqueness
case to a second endpoint.

For nonuniqueness, deform \(\log x\) inside a self-dual operator-monotone
family.  Self-duality keeps the equation traceless on the determinant-one
surface of real \(2\times2\) matrices, so the source's two-dimensional
Poincar\'e--Miranda mechanism survives.  One solution is calibrated exactly;
strict interval-certified signs on a rectangle force a second one.

\section{Duality and a second globally unique endpoint}

Call \(g\) \emph{globally unique} if \(\Gset_g(\omega;\mathbb A)\) is a
singleton for every \(n,m,\omega,\mathbb A\).

\begin{theorem}[Duality]\label{thm:duality}
For \(g\in\mathcal L\), define \(g^\vee(x)=-g(x^{-1})\).  Then
\(g^\vee\in\mathcal L\) and
\begin{equation}\label{eq:duality}
 \Gset_{g^\vee}(\omega;A_1,\ldots,A_m)
 =\{Y^{-1}:Y\in\Gset_g(\omega;A_1^{-1},\ldots,A_m^{-1})\}.
\end{equation}
Consequently, \(g\) is globally unique if and only if \(g^\vee\) is.
\end{theorem}

\begin{proof}
If \(0<A\le B\), then \(A^{-1}\ge B^{-1}\).  Operator monotonicity of
\(g\), followed by multiplication by \(-1\), shows
\(g^\vee(A)\le g^\vee(B)\).  Also \(g^\vee(1)=0\) and
\((g^\vee)'(1)=g'(1)=1\), so \(g^\vee\in\mathcal L\).

The metric geometric mean is self-dual:
\[
 (A_i\#X^{-1})^{-1}=A_i^{-1}\#X.
\]
Thus \(X\) satisfies the \(g^\vee\)-equation exactly when
\[
 0=-\sum_iw_i g(A_i^{-1}\#X).
\]
Putting \(Y=X^{-1}\) turns this into the \(g\)-equation for the inverse
data, proving \eqref{eq:duality}.  The uniqueness assertion follows by
applying the involution twice.
\end{proof}

\begin{corollary}\label{cor:endpoint}
The function \(g_-(x)=1-x^{-1}\) is globally unique, and its solution is
\[
 \Gset_{g_-}(\omega;\mathbb A)
 =\left\{\Omega(\omega;A_1^{-1},\ldots,A_m^{-1})^{-1}\right\},
\]
where \(\Omega\) denotes the Wasserstein mean.
\end{corollary}

\begin{proof}
The source recalls that \(g_+(x)=x-1\) is globally unique with solution
\(\Omega(\omega;\mathbb A)\).  Since \(g_+^\vee=g_-\), the claim is
Theorem~\ref{thm:duality}.
\end{proof}

\section{A non-logarithmic family with nonunique equations}

For \(0<s\le1\), set
\begin{equation}\label{eq:gs}
 g_s(x)=\frac{2}{s}\frac{x^s-1}{x^s+1},
 \qquad g_0(x):=\log x.
\end{equation}

\begin{proposition}\label{prop:gs}
For \(0<s\le1\), \(g_s\in\mathcal L\),
\(g_s(x^{-1})=-g_s(x)\), and \(g_s\to\log\) locally uniformly as
\(s\downarrow0\).
\end{proposition}

\begin{proof}
On the upper half-plane, the principal power \(z\mapsto z^s\) again has
positive imaginary part.  The real M\"obius map
\(z\mapsto(z-1)/(z+1)\) preserves the upper half-plane.  Hence \(g_s\) is a
Pick function, and L\"owner's analytic characterization makes it operator
monotone \cite[Chapter V]{Bhatia}.  Direct calculation gives
\(g_s(1)=0\), \(g_s'(1)=1\), and the self-duality identity.  The limit is
the scalar expansion \(x^s=e^{s\log x}\), uniform on compact subsets.
\end{proof}

Let
\[
 S=\begin{pmatrix}1&0\\0&-1\end{pmatrix},\qquad
 T=\begin{pmatrix}0&1\\1&0\end{pmatrix},\qquad X_0=e^{3S}.
\]
For a fixed sufficiently small \(s>0\), define
\begin{align*}
 a_s&=\frac2s\tanh\frac{3s}{2},
 &b_s&=\frac2s\tanh\frac{19s}{20},
 &r_s&=(a_s^2+b_s^2)^{1/2},\\
 q_s&=\frac2s\artanh\frac{s r_s}{2},
 &Z_1&=3S,
 &Z_2&=\frac{19}{10}T,\\
 Z_3&=-\frac{q_s}{r_s}(a_sS+b_sT),
 &B_i&=e^{Z_i},
 &A_i&=B_iX_0B_i.
\end{align*}
Here \(sr_s/2<1\) for small \(s\); in particular it holds at
\(s=1/100\), since \(a_s<3\) and \(b_s<19/10\).

\begin{theorem}[Explicit non-logarithmic nonuniqueness]\label{thm:nonunique}
Let \(s=1/100\), let \(g_s\) be given by \eqref{eq:gs}, take equal weights
\((1/3,1/3,1/3)\), and use the three matrices \(A_i\) above.  Then
\(\Gset_{g_s}(\omega;A_1,A_2,A_3)\) contains at least two distinct matrices.
Moreover, there is \(\varepsilon>0\) such that analogous examples are
nonunique for every \(0<s<\varepsilon\).
\end{theorem}

\begin{proof}
The relations \(S^2=T^2=I\) and \(ST=-TS\) give
\((a_sS+b_sT)^2=r_s^2I\).  If \(U\) is a self-adjoint involution, functional
calculus gives
\[
 g_s(e^{tU})=\frac2s\tanh\!\left(\frac{st}{2}\right)U.
\]
The definition of \(q_s\) therefore yields
\[
 g_s(B_1)=a_sS,\qquad g_s(B_2)=b_sT,\qquad
 g_s(B_3)=-a_sS-b_sT.
\]
Because \(A_i=B_iX_0B_i\), the Riccati characterization of \(\#\) shows
\(A_i\#X_0^{-1}=B_i\).  Hence \(X_0\) is an exact solution.

Now parametrize the real determinant-one positive matrices by
\[
 X(u,v)=e^{uS+vT}.
\]
Every \(A_i\) has determinant one, so does \(A_i\#X(u,v)^{-1}\).  Its two
eigenvalues are reciprocal; since \(g_s(x^{-1})=-g_s(x)\), its image under
\(g_s\) is traceless.  Consequently
\begin{equation}\label{eq:F}
 \frac13\sum_{i=1}^3g_s(A_i\#X(u,v)^{-1})
 =F_{1,s}(u,v)S+F_{2,s}(u,v)T
\end{equation}
for continuous real functions \(F_{1,s},F_{2,s}\).

Set
\[
 R=[1.61247432825,1.62347432825]
 \times[0.5194188906,0.5254188906].
\]
Outward-rounded interval arithmetic, subdividing each face into 8,000 exact
decimal intervals, certifies at \(s=1/100\):
\begin{align*}
 F_{1,s}(1.61247432825,v)&\ge1.5903\times10^{-4},\\
 F_{1,s}(1.62347432825,v)&\le-7.5949\times10^{-5},\\
 F_{2,s}(u,0.5194188906)&\ge1.1393\times10^{-4},\\
 F_{2,s}(u,0.5254188906)&\le-2.7880\times10^{-5}.
\end{align*}
The companion script records the complete enclosures.  The
Poincar\'e--Miranda theorem gives \((u_*,v_*)\in R\) with both functions in
\eqref{eq:F} zero.  Thus \(X(u_*,v_*)\) is a solution, and it differs from
\(X_0=X(3,0)\).

Finally, the calibrated matrices and all four face functions extend
continuously to \(s=0\).  The same interval script certifies strict face
signs there.  Uniform continuity on the four compact faces preserves those
signs for every sufficiently small \(s>0\), proving the last assertion.
\end{proof}

\section{Verification, limitations, and novelty}

The interval script uses the exact determinant-one identity
\[
 A\#B=\frac{A+B}{\sqrt{\det(A+B)}}\qquad
 (A,B\in\Pcone_2,\ \det A=\det B=1)
\]
and closed two-by-two functional calculus.  It reports zero failed
subintervals on every face for both \(s=0\) and \(s=1/100\).  The code is a
certificate only for the four scalar signs; the duality, operator
monotonicity, exact calibration, and topological conclusion are proved above.

The result does not give necessary and sufficient conditions for arbitrary
\(g\).  It shows that the uniqueness class is closed under duality, adds the
endpoint \(1-x^{-1}\), and proves that failure occupies a nontrivial
one-parameter neighborhood of the logarithmic mechanism.

The local run indexes were searched for arXiv:2605.16876 and the core
alternative-mean and generalized-Karcher terms.  Bounded web/arXiv searches
on 11 August 2026 used the exact title and question, \(1-x^{-1}\), and the
family \((x^s-1)/(x^s+1)\).  They found the May 2026 source and background
literature, but no later classification, no duality theorem for
\eqref{eq:gamma}, and no matching nonuniqueness family.  Novelty confidence is
moderate because the source is recent and the search was bounded.

Human review is recommended, focusing on Proposition~\ref{prop:gs}, the
factor in \(q_s\), and the interval implementation.

\begin{thebibliography}{9}
\bibitem{KM}
S. Kim and V. N. Mer,
\emph{A multivariable mean equation arising from the spectral geometric
mean}, arXiv:2605.16876 (2026).
\url{https://arxiv.org/abs/2605.16876}.

\bibitem{Bhatia}
R. Bhatia,
\emph{Matrix Analysis}, Graduate Texts in Mathematics 169,
Springer, 1997.
\end{thebibliography}

\end{document}

